\documentclass[11pt]{article}
\newcommand{\SheetCode}{Theory 06}
\usepackage[margin=25mm]{geometry}
\usepackage[T1]{fontenc}
\usepackage{lmodern}
\DeclareUnicodeCharacter{2605}{$\star$}
\usepackage{microtype}
\usepackage{amsmath,amssymb,mathtools,bm}
\usepackage{enumitem}
\usepackage{xcolor}
\usepackage{fancyhdr}
\usepackage{hyperref}
\usepackage[most]{tcolorbox}
\definecolor{agiBlue}{HTML}{173B57}
\definecolor{agiGold}{HTML}{B7791F}
\definecolor{agiLight}{HTML}{EFF6FA}
\definecolor{agiGreen}{HTML}{184D3B}
\hypersetup{colorlinks=true,linkcolor=agiBlue,urlcolor=agiBlue}
\setlength{\parindent}{0pt}
\setlength{\parskip}{0.55em}
\setlength{\headheight}{14pt}
\setlist{nosep,leftmargin=*}
\pagestyle{fancy}
\fancyhf{}
\fancyhead[L]{\small Part of the \href{https://github.com/mmikhasenko/GIModel.jl}{Agentic Gottried--Isgur} project}
\fancyhead[R]{\SheetCode}
\fancyfoot[C]{\thepage}
\newcommand{\dd}{\,\mathrm d}
\newcommand{\ket}[1]{\lvert #1\rangle}
\newcommand{\bra}[1]{\langle #1\rvert}
\newcommand{\braket}[2]{\langle #1\mid #2\rangle}
\newcommand{\expect}[1]{\left\langle #1\right\rangle}
\newcommand{\op}[1]{\widehat{#1}}
\newcommand{\GeV}{\,\mathrm{GeV}}
\newcommand{\R}{\mathbb{R}}
\newtcolorbox{goals}{colback=agiLight,colframe=agiBlue,title=Learning outcomes}
\newtcolorbox{reading}{colback=white,colframe=agiGold,title=Book reference,
  after skip=5mm}
\newtcolorbox{paperbridge}{colback=agiGreen!7,colframe=agiGreen,
  colbacktitle=agiGreen,coltitle=white,title=Connection to the GI paper}
\newtcolorbox{problem}[1]{colback=white,colframe=agiBlue,title={#1},
  before skip=3.5mm,after skip=1mm}
\newtcolorbox{solution}{colback=black!2,colframe=black!45,title=Solution}
\newtcolorbox{quiz}{colback=white,colframe=agiGold,title=Post-solution concept check}
\newtcolorbox{checkpoint}{colback=white,colframe=agiGold,title=Interpretation checkpoint}
\newcommand{\sheettitle}[3]{%
  \begin{center}
  {\Large\bfseries #1}\par
  \vspace{0.2em}{\color{agiBlue}#2}\par
  \vspace{0.4em}\small #3
  \end{center}\vspace{0.5em}}
\newcommand{\difficulty}[1]{\textbf{Difficulty:} #1}
\newcommand{\timebox}[1]{\textbf{Suggested time:} #1}
\newcommand{\answer}{\par\smallskip\color{black!50}\textbf{Answer:} }

\begin{document}
\sheettitle{Sheet 6: Resolving the multiplets}{Contact, tensor, and spin--orbit structure}{Prerequisite: Sheets 1--5. \difficulty{★★/★★★} \quad \timebox{3--4 hours}}

\begin{goals}
Connect radial interactions to angular operators; derive contact selection and
singlet/triplet patterns; use angular sum rules; distinguish vector and scalar
spin--orbit terms; understand why spin terms may be included directly in a
fixed-channel diagonalization instead of only to first order.
\end{goals}
\begin{reading}
Selected reading in Landau--Lifshitz, \emph{Quantum Mechanics: Non-Relativistic
Theory}: Ch. VI, ``Perturbation Theory,'' Ch. VIII, ``Spin,'' and Ch. XIV,
``Addition of Angular Momenta.'' This is optional background rather than the
sequence or notation that the sheet is required to follow.
\end{reading}


Throughout this sheet, a \emph{fixed channel} means that $L$, $S$, and $J$ are
held fixed while the radial Hamiltonian is solved. Mixing between different
spectroscopic channels is postponed to Sheet 7.

\begin{problem}{Problem 1 ★ — Prove the unsmeared contact selection rule}
Show that a regular $L>0$ wave has zero expectation of $\delta^3(\bm r)$, while
an $S$ wave gives $|\Psi(0)|^2$.

\textbf{Hint:} Use the regular branch $R_L(r)\sim r^L$ from Sheet 1 and
the defining identity $\int f(\bm r)\delta^3(\bm r)\dd^3r=f(0)$
for a continuous density $f$.
\end{problem}
\begin{solution}
$\int|\Psi|^2\delta^3(\bm r)\dd^3r=|\Psi(0)|^2$. Regularity gives
$R_L\sim r^L$, so the value vanishes for $L>0$ and may be finite for $L=0$.
Gaussian smearing broadens this exact selection rule into a strongly
short-distance-weighted expectation.
\end{solution}
\begin{quiz}
Why may a smeared contact term give a small nonzero $P$-wave contribution?
\answer it samples a finite ball rather than only the point $r=0$.
\end{quiz}

\begin{problem}{Problem 2 ★ — Derive the hyperfine ratio}
Two spin-$1/2$ constituents share the same normalized radial wave.
Let $H_c=C\bm S_1\cdot\bm S_2$, where the fixed real energy $C$ already
includes the radial matrix element. Compute the first-order singlet and
triplet shifts $\Delta E_0,\Delta E_1$, their ratio for $C\ne0$, and the
centroid shift $(\Delta E_0+3\Delta E_1)/4$.

\textbf{Hint:} Use the spin expectations from Sheet 2, Problem 3;
the weights are $2S+1$.
\end{problem}
\begin{solution}
$\Delta E_{S=0}=-3C/4$ and $\Delta E_{S=1}=C/4$,
so $\Delta E_0/\Delta E_1=-3$ for $C\ne0$. Weighting by spin
degeneracies gives $1(-3C/4)+3(C/4)=0$. Radial distortion in a full
diagonalization can modify the simple ratio.
\end{solution}
\begin{quiz}
If the reported contact shifts do not have exactly ratio $-3$, must the code be
wrong? \answer not if singlet and triplet Hamiltonians were diagonalized
separately and acquired different radial waves.
\end{quiz}

\begin{problem}{Problem 3 ★★ — Separate vector and scalar spin--orbit physics}
For equal constituent masses $m$, take as input the leading
nonrelativistic vector/scalar spin--orbit reduction
\[
H_{LS}=\frac{3V_V'(r)-V_S'(r)}{2m^2r}\bm L\cdot\bm S.
\]
Use a Lorentz-vector potential $V_V=-\kappa/r$ and a Lorentz-scalar potential
$V_S=br+c$, with $\kappa,b>0$. Define $H_{LS}=C_{LS}(r)\bm L\cdot\bm S$.
Find the two contributions to $C_{LS}(r)$
and compare their signs and radial dependence. The reduction formula is
supplied; deriving it from relativistic dynamics is not required. Smearing
and momentum factors are omitted in this exercise.
\end{problem}
\begin{solution}
Differentiating gives $C_{LS}^{(V)}=3\kappa/(2m^2r^3)>0$ and
$C_{LS}^{(S)}=-b/(2m^2r)<0$.
The color-magnetic vector interaction generates a direct spin--orbit coupling,
whereas scalar confinement contributes through Thomas precession of the
accelerating rest frame. Their radial derivatives enter with different signs.
The observed splitting therefore tests both the Lorentz structure and the
short-distance kernel, not merely one coefficient.
\end{solution}
\begin{quiz}
Can the central spectrum alone determine the Lorentz nature of confinement?
\answer weakly at best; spin--orbit splittings carry the discriminating
information.
\end{quiz}

\begin{problem}{Problem 4 ★★ — Use triplet centroid identities}
For $^3P_J$, the conventional tensor expectations are
$t_J=(-4,2,-2/5)$ for $J=(0,1,2)$. Verify
$\sum_J(2J+1)t_J=0$ and combine this with the Sheet-2 spin--orbit sum rule.

\textbf{Hint:} Assume a common radial wave and first-order shifts
$\Delta E_J=A\expect{\bm L\cdot\bm S}_J+Bt_J$, with $A,B$ independent
of $J$. Determine the shift of $\overline E=\sum_J(2J+1)E_J/9$.
Here $t_J$ uses $S_{12}=12(\bm S_1\cdot\widehat{\bm r})
(\bm S_2\cdot\widehat{\bm r})-4\bm S_1\cdot\bm S_2$.
\end{problem}
\begin{solution}
$1(-4)+3(2)+5(-2/5)=0$. Since the spin--orbit weighted sum also vanishes,
first-order tensor and spin--orbit corrections leave the degeneracy-weighted
$^3P$ centroid unchanged. This is a useful algebraic check because it does not
depend on the radial wavefunction.
\end{solution}
\begin{quiz}
What kind of error survives this centroid check? \answer an incorrect radial
normalization or common multiplicative error in the radial interaction can
scale all shifts while preserving the angular identity.
\end{quiz}

\begin{problem}{Problem 5 ★★★ — Show why tensor forces mix $L$ and $L+2$}
Treat the tensor interaction as the scalar coupling of a rank-2 spatial tensor
and a rank-2 spin tensor. Use the triangle rule and parity to derive
$\Delta L=0,\pm2$ for nonzero matrix elements.

\textbf{Hint:} A spatial tensor of rank $k$ has a nonzero angular matrix
element only if $|L-L'|\le k\le L+L'$. For $Y_{2q}$ parity requires
$L+L'+2$ even. State these as necessary conditions, including the restriction
on $L+L'$; $\Delta L=0$ alone does not permit an $L=L'=0$ matrix element.
The full scalar interaction conserves $J$.
\end{problem}
\begin{solution}
A rank-2 spherical harmonic permits $|L-L'|\le2\le L+L'$. Parity conservation
requires $L$ and $L'$ to have the same parity, eliminating $\Delta L=\pm1$.
Thus diagonal and $\pm2$ matrix elements remain. For example, $^3S_1$ can mix
with $^3D_1$.
\end{solution}
\begin{quiz}
Why is the tensor operator zero in a spin singlet? \answer two spins coupled to
$S=0$ cannot support a rank-2 spin tensor.
\end{quiz}

\begin{problem}{Problem 6 ★★★ — Compare perturbation and fixed-channel solution}
In one fixed $(L,S,J)$ channel, let $H(\lambda)=H_0+\lambda V_s$,
where $V_s$ is Hermitian and independent of the dimensionless parameter
$\lambda$. A normalized, nondegenerate state satisfies
$H_0\ket n=E_n^{(0)}\ket n$.
Derive the coefficient $E_n^{(1)}$ in
$E_n(\lambda)=E_n^{(0)}+\lambda E_n^{(1)}+O(\lambda^2)$.
Then explain what a direct diagonalization of $H(1)$ can change in both
the energy and the radial wave used for later mixing.

\textbf{Hint:} Differentiate the eigenvalue equation at $\lambda=0$
and multiply on the left by $\bra n$.
\end{problem}
\begin{solution}
For a nondegenerate normalized eigenstate of $H_0$,
$E_n'(0)=\bra nV_s\ket n$. Diagonalizing $H(1)$ includes the effect of $V_s$
to all orders within the chosen finite basis, lets the radial wave respond to
contact and fine structure, and makes
later cross-channel elements use the physically distorted fixed-channel waves.
It also makes the reported operator contributions expectations in one common
eigenstate whose sum equals the eigenvalue.
\end{solution}
\begin{quiz}
Why should reported ``central + contact + tensor + spin--orbit'' pieces use the
same wave? \answer otherwise they do not decompose one Hamiltonian expectation
and need not add to the stated mass.
\end{quiz}

\begin{paperbridge}
GI Eq. (4) contains the contact and tensor hyperfine interactions; Eqs. (5)--(7)
separate color-magnetic and Thomas-precession spin--orbit terms. Their smeared,
momentum-dependent forms appear in Appendix-A Eqs. (A15)--(A16). The first
stage described after Eq. (14) diagonalizes the diagonal pieces within fixed
$\ket{jm;ls}$ sectors.
\end{paperbridge}
\end{document}
